\chapter{Thermodynamics of Combustion and Chemical Equilibrium}\label{ch:thermo}

This chapter develops the two coarsest levels of analysis introduced in Section~\ref{ch:intro}: the pure enthalpy/heat balance, which yields heating values and the adiabatic flame temperature, and chemical equilibrium, which corrects those estimates for dissociation and predicts how products and temperature respond to pressure and preheat. Throughout we treat the gas as a mixture of thermally perfect gases and the process as either isobaric (the relevant idealisation for gas-turbine combustors) or isochoric~\cite{lecture2,Langella_LectureNotes}.

\section{Perfect gases, the laws of thermodynamics, and why combustors run at constant pressure}

A thermally perfect gas obeys $p = \rho R T = (\rho/W) R_0 T$, with $R_0=\SI{8.314}{J\,mol^{-1}\,K^{-1}}$ the universal gas constant and $R=R_0/W$ the specific gas constant~\cite{Langella_LectureNotes}. The first law for a closed system undergoing a quasi-static change is $\dd u = \delta q - \delta w = \delta q - p\,\dd v$, where lower-case symbols are per unit mass. Two process constraints are of interest. At \emph{constant volume} the work term vanishes and $\delta q = \dd u = c_v\,\dd T$, so $c_v$ measures how internal energy is stored. At \emph{constant pressure} it is convenient to introduce the enthalpy $h \equiv u + pv$; then $\dd h = \dd u + p\,\dd v + v\,\dd p$ and, since $\dd p = 0$, $\delta q = \dd h = c_p\,\dd T$~\cite{Langella_LectureNotes,warnatz2006combustion}.

This last identity is the reason combustion is overwhelmingly analysed at constant pressure: in a nearly isobaric flow (a deflagration, hence essentially every air-breathing combustor) the heat exchanged equals an enthalpy change directly, $\delta q = \dd h$, \emph{without any reference to the chemical pathway}, because $h$ is a state function~\cite{Langella_LectureNotes}. For a perfect gas $c_p = c_v + R$ and $\gamma = c_p/c_v$. The lecture poses the conceptual question of the ``advantage of an isobaric over an isochoric process'': at constant pressure part of the added heat goes into expansion work rather than into raising the temperature, so $c_p>c_v$, and a given heat addition produces a smaller temperature (and entropy) rise --- but the bookkeeping via enthalpy is exact and path-independent, which is what makes the isobaric idealisation so powerful for design~\cite{lecture1,transcript_tutorial}.

\section{Heat capacity and why it varies}

Heat capacity is the heat required to raise the temperature of a system by one kelvin, $C=\delta Q/\dd T$; per unit mass $c=\delta q/\dd T$ and per mole $\bar C = W c$~\cite{Langella_LectureNotes}. Because heat is not a state function, $c$ is only well defined once the process is fixed, giving the two cases $c_p$ and $c_v$ above. The variation of $c$ with temperature is governed by the molecular \emph{degrees of freedom} over which energy is equipartitioned~\cite{lecture1,warnatz2006combustion}. A monatomic gas stores energy only in three translational modes, giving $\bar C_v=\tfrac32 R_0$ independent of temperature. A diatomic molecule activates two rotational modes at moderate temperatures ($\bar C_v\to\tfrac52 R_0$) and, at high temperature, vibrational modes ($\bar C_v\to\tfrac72 R_0$); these modes are quantised and only become accessible (``unlock'') when $k_BT$ is comparable to their energy spacing. Polyatomic molecules have still more rotational and vibrational modes, and they activate at lower temperatures, so their heat capacity is larger and rises faster with $T$~\cite{lecture1,warnatz2006combustion}.

Since temperature measures only the chaotic translational motion, when new internal modes open up an increasing share of any added heat is stored internally rather than as translational kinetic energy: more heat is needed per kelvin and $c$ increases monotonically with $T$~\cite{lecture1}. Heat capacity also depends on \emph{composition}, since for a mixture $c_{p,\mathrm{mix}}=\sum_k Y_k c_{p,k}$; as combustion converts reactants into triatomic \ch{CO_2} and \ch{H_2O}, which have larger $c_p$, the mixture heat capacity changes even at fixed temperature~\cite{warnatz2006combustion,glassman2014combustion}. This double dependence $c_p(T,Y_k)$ is exactly why accurate calculations tabulate $c_{p,k}(T)$ as polynomial fits --- the NASA seven-coefficient polynomials used by tools such as Cantera --- rather than constant values, and it is the root cause of the parameter sensitivities of the adiabatic flame temperature discussed below~\cite{warnatz2006combustion,cantera,transcript_tutorial}.

\begin{figure}[htbp]
\centering
\includegraphics[width=0.62\linewidth]{fig_heat_capacity}
\caption{Variation of the molar heat capacity with temperature, showing the successive activation of translational, rotational and vibrational degrees of freedom (limits $\tfrac32 R_0,\ \tfrac52 R_0,\ \tfrac72 R_0$)~\cite{lecture1}.}
\label{fig:fig_heat_capacity}
\end{figure}

\section{The types of enthalpy and the heat of reaction}

Because no absolute enthalpy can be defined, all values are referred to a convention~\cite{lecture1,Langella_LectureNotes}. The \textbf{standard enthalpy of formation} $\Delta h_f^0$ of a species is the enthalpy change when one unit amount of it is formed from its constituent elements with all substances in their standard states ($p=\SI{1}{bar}$, $T_0=\SI{298.15}{K}$, chosen over \SI{0}{K} for experimental practicality)~\cite{lecture1}. Elements in their most stable form (\ch{O_2}, \ch{N_2}, \ch{H_2}, graphite) are assigned $\Delta h_f^0=0$ by convention. Formation enthalpies are measured calorimetrically (the bomb calorimeter) and combined using \textbf{Hess's law} --- the algebraic sum of the heats evolved in separate stages equals the heat of the direct reaction, a direct consequence of $h$ being a state function --- to obtain values for species whose formation reaction cannot be run directly~\cite{lecture1,warnatz2006combustion,maas_combustion_lecture}.

The \textbf{absolute (total) enthalpy} of a species is the sum of its chemical and thermal (sensible) parts,
\begin{equation}
h_k(T)=\Delta h_{f,k}^0+\int_{T_0}^{T}c_{p,k}(T')\,\dd T',\qquad
h_{\mathrm{mix}}=\sum_k Y_k h_k ,
\label{eq:abs-enthalpy}
\end{equation}
the integral being the sensible enthalpy~\cite{lecture1,warnatz2006combustion}. Phase changes add the latent enthalpies of vaporisation and fusion, relevant for liquid fuels and for the state of product water. The \textbf{enthalpy (heat) of reaction} $\Delta H_R$ is the heat that must be removed (exothermic) or added (endothermic) to return the products to the reactant temperature; the constant-temperature requirement is what makes the definition unique. From Hess's law,
\begin{equation}
(\Delta Q_R)_p=\Delta H_R=\Big(\textstyle\sum_{\mathrm{prod}}\nu_i W_i\,\Delta h_{f,i}^0-\sum_{\mathrm{reac}}\nu_i W_i\,\Delta h_{f,i}^0\Big)+\int_{T_0}^{T}\Delta c_p(T')\,\dd T',
\label{eq:heat-reaction}
\end{equation}
where the last term (\textbf{Kirchhoff's law}) corrects for reactants not at standard temperature and is usually small~\cite{lecture1,Langella_LectureNotes,glassman2014combustion}. All fuels have negative heats of reaction (combustion is exothermic), even though some, like \ch{C_2H_2}, have \emph{positive} enthalpies of formation~\cite{lecture1}. Because product water may be counted as liquid or vapour, one distinguishes the \textbf{higher heating value} (HHV, water condensed, latent heat recovered) from the \textbf{lower heating value} (LHV, water remains vapour); the LHV is the relevant one for combustors, where residence times are far too short for condensation~\cite{Langella_LectureNotes,glassman2014combustion}. A useful set of facts from the lecture: the heats of reaction of most hydrocarbons cluster at \SIrange{40}{50}{MJ\,kg^{-1}}, their stoichiometric air-fuel ratios between about 13:1 and 17:1, and alcohols have lower mass-specific heats of reaction because their bound oxygen contributes no energy~\cite{lecture1}.

\section{The adiabatic flame temperature}

The \textbf{adiabatic flame temperature} $\Tad$ is the product temperature when combustion proceeds with no heat loss. For the isobaric case it follows from enthalpy conservation, $h_{\mathrm{prod}}(\Tad)=h_{\mathrm{reac}}(T_{\mathrm{reac}})$, i.e.\ all the chemical energy released is used to heat the products~\cite{lecture1,lecture2}. Writing this out with Eq.~\eqref{eq:abs-enthalpy} and assuming complete, one-step, infinitely fast chemistry gives, after relating the species changes through the fuel via $\Delta Y_F/(\nu_F W_F)=\Delta Y_k/(\nu_k W_k)$, the explicit estimate~\cite{lecture2}
\begin{equation}
\Tad = T_0 + \frac{Y_{F,r}\,\Delta H_R}{c_{p,p}\,\nu_F W_F} + \frac{c_{p,r}}{c_{p,p}}\,(T_r-T_0)\qquad(\phi\le1),
\label{eq:Tad}
\end{equation}
with the fuel mass fraction replaced by the oxidiser one for $\phi>1$ (the deficient reactant limits the heat release)~\cite{lecture2}. The equivalence ratio is
$\phi=(\dot m_F/\dot m_{Ox})/(\dot m_F/\dot m_{Ox})_{st}=(Y_F/Y_{Ox})/(Y_F/Y_{Ox})_{st}$, equal to unity at stoichiometric, below for lean and above for rich mixtures~\cite{lecture2}. Equation~\eqref{eq:Tad} exposes the two governing factors: the energy released \emph{per unit mass of mixture} and the heat capacity of the products. Representative stoichiometric values in air are \SI{2263}{K} for \ch{CH_4}, \SI{2289}{K} for \ch{C_2H_6}, \SI{2290}{K} for \ch{C_3H_8} and \SI{2439}{K} for \ch{C_2H_4}~\cite{lecture2}.

\begin{figure}[htbp]
\centering
\includegraphics[width=0.6\linewidth]{fig_tad_phi}
\caption{Adiabatic flame temperature versus equivalence ratio, peaking slightly rich of stoichiometric~\cite{lecture2}.}
\label{fig:fig_tad_phi}
\end{figure}

\paragraph{Equivalence ratio.} $\Tad$ peaks near stoichiometric, where neither excess air nor excess fuel acts as inert ballast absorbing heat, but the peak lies slightly \emph{rich} (typically $\phi\approx1.05$--$1.1$)~\cite{lecture1,lecture2}. There are two reasons: on the rich side the products contain more \ch{CO} and \ch{H_2} and less \ch{CO_2}/\ch{H_2O}, lowering the mean product $c_p$; and dissociation losses are reduced because less \ch{CO_2} is formed to begin with. Far from stoichiometric, $\Tad$ falls steeply on both sides~\cite{lecture2,glassman2014combustion}.

\paragraph{Pressure.} The complete-combustion analysis of Eq.~\eqref{eq:Tad} shows no pressure dependence, yet measured and equilibrium temperatures \emph{increase} with pressure. The reason is dissociation, addressed in the next section through Le~Chatelier's principle~\cite{lecture2}.

\paragraph{Inlet temperature.} Preheating the reactants raises $\Tad$, but by \emph{less} than the preheat increment: the product heat capacity exceeds that of the reactants ($c_{p,r}/c_{p,p}<1$ in Eq.~\eqref{eq:Tad}), and the higher final temperature enhances endothermic dissociation, which absorbs part of the added sensible energy~\cite{lecture2,transcript_tutorial}.

\paragraph{Process and fuel.} For an \emph{isochoric} process $\Tad$ is higher than for the isobaric case, because no energy is spent on expansion work ($\dd u=0$ instead of $\dd h=0$, and $c_v<c_p$)~\cite{lecture1}. Removing the nitrogen ballast has a dramatic effect: \ch{H_2}/air gives roughly \SI{2380}{K} but \ch{H_2}/\ch{O_2} about \SI{3083}{K}~\cite{maas_combustion_lecture}.

\section{Chemical equilibrium and free energy}

The complete-combustion picture fails at two points: it ignores the \emph{dissociation} of products (e.g.\ $\ch{CO_2}\rightleftharpoons\ch{CO}+\tfrac12\ch{O_2}$, $\ch{H_2O}\rightleftharpoons\ch{H_2}+\tfrac12\ch{O_2}$), which is endothermic and becomes important at the high temperatures near stoichiometric, and it cannot predict minor species at all~\cite{lecture2}. The equilibrium description removes the first limitation by assuming the chemistry is infinitely fast, so that the composition is the one that the reactions, run forward and backward, no longer change.

\subsection{From the second law to the Gibbs function}

One could find equilibrium directly from kinetics: the rate of progress of reaction $j$ is $Q_j=K_{f,j}\prod_k X_k^{\nu'_{k,j}}-K_{r,j}\prod_k X_k^{\nu''_{k,j}}$, and setting $Q_j=0$ gives the concentration-based equilibrium constant $K_{c,j}=K_{f,j}/K_{r,j}=\prod_k X_{k,e}^{\nu_{k,j}}$ (the law of mass action)~\cite{lecture2}. But this requires the full rate data. Thermodynamics offers a shortcut. The second law states that a spontaneous process increases the entropy of the universe, $\Delta S_{\mathrm{system}}+\Delta S_{\mathrm{surr}}>0$, but $\Delta S_{\mathrm{surr}}$ is inconvenient. At constant temperature and pressure the heat passed to the surroundings is $\Delta Q_{\mathrm{surr}}=-\Delta H_{\mathrm{system}}$ (first law, isobaric), and $\Delta S_{\mathrm{surr}}=\Delta Q_{\mathrm{surr}}/T=-\Delta H_{\mathrm{system}}/T$; substituting and rearranging expresses the criterion in system properties alone~\cite{lecture2}:
\begin{equation}
\Delta G\equiv\Delta H-T\Delta S<0,\qquad G=H-TS .
\label{eq:gibbs}
\end{equation}
The \textbf{Gibbs free energy} $G$ is a state function representing the ``available'' energy: the part $TS$ is tied up in entropy and unavailable for work, so $\Delta G$ is the excess energy liberated out of equilibrium~\cite{lecture2}. For an isochoric process the \textbf{Helmholtz free energy} $F=U-TS$ takes its place. If $\Delta G<0$ the reaction is exergonic and proceeds forward; if $\Delta G>0$ it proceeds in reverse; if $\Delta G=0$ the system is in (dynamic) equilibrium and $G$ is at a minimum~\cite{lecture2,transcript_tutorial}. Spontaneity therefore depends on temperature and is \emph{not} the same as exothermicity: a reaction with $\Delta H<0,\Delta S>0$ is spontaneous at all $T$; one with $\Delta H>0,\Delta S>0$ (such as the dissociation $\ch{N_2}\rightarrow\ch{N}+\ch{N}$) only at high $T$ where $T\Delta S$ dominates; one with $\Delta H<0,\Delta S<0$ only at low $T$; and one with $\Delta H>0,\Delta S<0$ never~\cite{lecture2}. This is precisely why dissociation reactions switch on at flame temperatures.

\subsection{The law of mass action and the equilibrium constant}

To extract the equilibrium composition, take $G=H-TS$ for an ideal gas in a closed system and differentiate, using the Gibbs relation $\dd U=T\,\dd S-p\,\dd V$, to obtain $\dd G=V\,\dd p-S\,\dd T$. At constant temperature, $G-G^0=\int_{p^0}^{p} nR_0T\,\dd p'/p' = nR_0T\ln(p/p^0)$, so for the reaction $\Delta G_R=\Delta G_R^0+R_0T\ln\prod_k (p_k/p^0)^{\nu_k}$~\cite{lecture2}. Setting $\Delta G_R=0$ at equilibrium yields the \textbf{law of mass action}
\begin{equation}
\prod_{k=1}^{N}\left(\frac{p_k}{p^0}\right)^{\nu_k}=K_p=\exp\!\left(-\frac{\Delta G_R^0}{R_0 T}\right),\qquad K_p=K_c\,(R_0T)^{\Delta n},
\label{eq:Kp}
\end{equation}
with $\Delta n=\sum_k\nu_k$ the change in the number of moles and $p_k=pX_k$~\cite{lecture2,warnatz2006combustion}. Together with element conservation and the enthalpy balance of Eq.~\eqref{eq:abs-enthalpy}, Eq.~\eqref{eq:Kp} fixes the equilibrium products and, self-consistently, the adiabatic flame temperature~\cite{lecture2,glassman2014combustion}.

\subsection{Le Chatelier's principle: pressure and preheat revisited}

Equation~\eqref{eq:Kp} now explains the pressure effect. Dissociation reactions increase the number of moles ($\Delta n>0$); since $\prod_k(p_k/p^0)^{\nu_k}=K_p(T)$ depends only on temperature, raising the total pressure forces the dissociation-product mole fractions down --- \textbf{Le~Chatelier's principle}: the system shifts toward the side with fewer moles~\cite{lecture2,transcript_tutorial}. Because dissociation is endothermic, suppressing it leaves more energy as sensible heat, so $\Tad$ \emph{rises with pressure}; the effect saturates at high pressure and is strongest near stoichiometric, where temperatures and hence dissociation are largest~\cite{lecture2,glassman2014combustion}. Conversely, raising the temperature shifts equilibrium toward the endothermic (more dissociated) side, which is the microscopic reason the preheat gain in $\Tad$ is less than one-for-one~\cite{transcript_tutorial}.

\begin{keybox}
\keyhead{Dissociation, pressure and temperature}
Higher pressure suppresses (mole-increasing) dissociation $\Rightarrow$ more complete combustion $\Rightarrow$ higher $\Tad$. Higher temperature promotes dissociation $\Rightarrow$ buffers further temperature rise. Both follow from $\prod_k(p_k/p^0)^{\nu_k}=K_p(T)$ and Le~Chatelier's principle~\cite{lecture2}.
\end{keybox}

\subsection{When equilibrium is and is not enough}

Equilibrium assumes infinitely fast chemistry, equivalent to infinite available time. In reality a chemical time scale competes with the residence time. For the \emph{major} products the equilibrium approximation is good once the gas has had enough time --- by the combustor exit the major species are close to equilibrium~\cite{transcript_tutorial}. For the \emph{minor} species \ch{CO} and \ch{NO_x} it fails: their formation/consumption time scales are comparable to or longer than the residence time. Equilibrium over-predicts the conversion of \ch{CO} to \ch{CO_2} (so measured \ch{CO} exceeds the equilibrium value), and equilibrium \ch{NO_x} is generally not representative of measured values, because \ch{NO_x} forms slowly through radical pathways~\cite{transcript_tutorial}. This is exactly the gap that finite-rate chemical kinetics, the subject of Chapter~\ref{ch:kinetics}, must fill, and it is why the course assignment contrasts equilibrium predictions with finite-rate flame computations for these species~\cite{transcript_tutorial}.
